% Amendment A2, source: results/phase2/a2/analysis/real_final/{table_a2.md,a2_paired.csv}; tiles: data/real_tiles/hirise_tiles.json
\textbf{Real-terrain transfer (pre-registered amendment A2).} To test whether the conclusions depend on the procedural
generator, we built $32\times32$ maps from three public HiRISE digital terrain models at rover sites: Victoria crater
(Opportunity), Jezero crater (Perseverance) and Gale crater (Curiosity). Elevations were averaged to 2\,m cells. Each
cell was classified by a rule fixed before the data was inspected: obstacle if slope $>25^\circ$ or local relief above
the site's 97th percentile; rock if slope is 15--25$^\circ$; sand if slope $<5^\circ$ and relief below the 20th
percentile; normal otherwise. We kept fully valid, non-overlapping tiles that pass the simulator's own acceptance rule,
with at most 300 per site, split half for debugging and half for confirmation (450 confirmatory tiles; Fig.~\ref{fig:real}).
Physics, models and thresholds are the in-distribution ones, so this is an evaluation-only test, frozen before any run.
The class rule encodes \emph{terrain structure} from real elevation data, not measured traversability.

Real terrain is flatter than our procedural world (3--7\% obstacle cells against 15\%), and every method does better on
it (Table~\ref{tab:real}). The qualitative conclusions hold. \ours{} again has the lowest unsafe-action exposure of all
baselines (0.019 attempted unsafe moves per episode, against 0.040 for A*+monitor and 0.23 for DQN+shield), and the
monitor-only variant reaches the same level (0.014), so the monitor again carries the safety gain. Unlike on procedural
maps, the learned system no longer trails the classical planners in success. It matches plain A* (0.996 vs.\ 0.996,
difference $+0.0002$, n.s.) with 81\% fewer attempted unsafe moves, 0.45 fewer collisions and 3.1 fewer energy units per
episode. It also beats A*+monitor on success ($+0.5$ points), unsafe moves ($-0.021$), collisions ($-0.38$) and energy
($-4.3$ units); all primary comparisons are significant after Holm correction. The full systems intervene on only 1.5\%
(action-conditioned OOD) and 3.3\% (raw OOD) of decisions. The action-conditioned variant does not change success and
lowers unsafe moves only marginally ($-0.004$ per episode).

\begin{figure}[t]
\centering
\includegraphics[width=\columnwidth]{fig10_real_tiles}
\caption{Amendment A2: example confirmatory tiles derived from HiRISE DTMs (2\,m cells; colours as in Fig.~2) and
the terrain composition per site compared with the procedural ID generator.}
\label{fig:real}
\end{figure}

\begin{table}[t]
\centering
\caption{Real-terrain transfer (450 HiRISE tiles, ID physics; mean over 10 seeds). Source: \texttt{table\_a2.md}.}
\label{tab:real}
\footnotesize
\setlength{\tabcolsep}{3pt}
\begin{tabular}{@{}lcccc@{}}
\toprule
Method & Success & Unsafe/ep. & Collisions/ep. & Energy \\ \midrule
A*                 & .996 & .102 & .569 & 79.7 \\
A*+monitor         & .991 & .040 & .496 & 80.8 \\
DQN                & .965 & 3.20 & 3.31 & 84.0 \\
DQN+shield         & .972 & .231 & .398 & 83.1 \\
PPO+shield         & .962 & .508 & .631 & 86.5 \\
Ensemble           & .980 & 1.90 & 1.98 & 78.8 \\
Unc.-gated         & .992 & 1.07 & 1.16 & 77.7 \\
\ours{}            & \textbf{.996} & .019 & .118 & \textbf{76.6} \\
\ours{}, action-OOD & .994 & .015 & .117 & 76.9 \\
Monitor only       & .987 & \textbf{.014} & \textbf{.115} & 77.2 \\ \bottomrule
\end{tabular}
\end{table}
